\fiche{Concordance du programme}{Des attendus aux chapitres et aux exercices}
\phantomsection\label{concordance}\pdfbookmark[0]{Programme et sources}{programme}
Référence : programme de l'enseignement optionnel de mathématiques expertes de Terminale générale, BO spécial n° 8 du 25 juillet 2019, arrêté du 19 juillet 2019, NOR MENE1921264A.
La page Éduscol consultée le 23 octobre 2026 conserve ce programme pour l'option. Les nouveaux programmes publiés en avril 2026 pour d'autres enseignements ne doivent pas lui être substitués.
Les regroupements ci-dessous sont éditoriaux ; les exemples de problèmes possibles du programme permettent un choix pédagogique et ne sont pas tous des obligations.
\par\vspace{3mm}
\noindent\textbf{C01} · Ensemble complexe, parties réelle et imaginaire, opérations et équation $az=b$. \hyperref[chap-1]{Chapitre 1}, exercices 1, 2, 5.\par\vspace{1.7mm}
\noindent\textbf{C02} · Conjugaison, inverse, équations simples avec $z$ et son conjugué. \hyperref[chap-1]{Chapitre 1}, exercices 3, 4, 6, 7, 9, 11.\par\vspace{1.7mm}
\noindent\textbf{C03} · Conjugué du produit, de l’inverse et des puissances ; formule du binôme et sa preuve. \hyperref[chap-1]{Chapitre 1}, exercices 8, 10, 12.\par\vspace{1.7mm}
\noindent\textbf{C04} · Affixes de points/vecteurs, image du conjugué et représentation. \hyperref[chap-2]{Chapitre 2}, exercices 1, 5.\par\vspace{1.7mm}
\noindent\textbf{C05} · Module, distance, $|z|^2=z\bar z$, produit, inverse et puissance. \hyperref[chap-2]{Chapitre 2}, exercices 2, 6, 9.\par\vspace{1.7mm}
\noindent\textbf{C06} · Cercle unité et stabilité ; arguments et forme trigonométrique. \hyperref[chap-2]{Chapitre 2}, exercices 3, 4, 6, 8.\par\vspace{1.7mm}
\noindent\textbf{C07} · Addition et duplication trigonométriques ; démonstration par produit scalaire. \hyperref[chap-3]{Chapitre 3}, exercices 4, 5, 11.\par\vspace{1.7mm}
\noindent\textbf{C08} · Exponentielle imaginaire, relation fonctionnelle, changement des trois formes. \hyperref[chap-3]{Chapitre 3}, exercices 1, 2, 3, 8.\par\vspace{1.7mm}
\noindent\textbf{C09} · Euler et Moivre : transformations, intégrales, suites et puissances. \hyperref[chap-3]{Chapitre 3}, exercices 6, 7, 9, 10, 12.\par\vspace{1.7mm}
\noindent\textbf{C10} · Second degré réel dans $\C$, racine connue d’un troisième degré. \hyperref[chap-4]{Chapitre 4}, exercices 1, 2, 4, 5, 6.\par\vspace{1.7mm}
\noindent\textbf{C11} · Factorisations $z^n-a^n$ et $P(z)$ par $z-a$ ; borne du nombre de racines et preuves. \hyperref[chap-4]{Chapitre 4}, exercices 3, 7, 8, 9, 10, 12.\par\vspace{1.7mm}
\noindent\textbf{C12} · Module et argument d’un quotient de vecteurs, alignement, orthogonalité et lieux. \hyperref[chap-5]{Chapitre 5}, exercices 1, 2, 5, 6, 8.\par\vspace{1.7mm}
\noindent\textbf{C13} · Racines de l’unité, preuve, cas $2,3,4$, polygones réguliers. \hyperref[chap-5]{Chapitre 5}, exercices 3, 4, 7, 12.\par\vspace{1.7mm}
\noindent\textbf{A01} · Divisibilité dans $\Z$, combinaisons et diviseurs. \hyperref[chap-6]{Chapitre 6}, exercices 1, 5, 10, 12.\par\vspace{1.7mm}
\noindent\textbf{A02} · Division euclidienne d’un relatif par un entier strictement positif. \hyperref[chap-6]{Chapitre 6}, exercices 2, 3.\par\vspace{1.7mm}
\noindent\textbf{A03} · Congruences et compatibilité avec opérations ; tests de divisibilité. \hyperref[chap-6]{Chapitre 6}, exercices 4, 6, 7, 8, 9, 11.\par\vspace{1.7mm}
\noindent\textbf{A04} · PGCD, Euclide, couples premiers entre eux et Bézout constructif. \hyperref[chap-7]{Chapitre 7}, exercices 1, 2, 4, 5.\par\vspace{1.7mm}
\noindent\textbf{A05} · Gauss : démonstration et applications. \hyperref[chap-7]{Chapitre 7}, exercices 6, 12.\par\vspace{1.7mm}
\noindent\textbf{A06} · Inverse modulaire et résolution de $ax\equiv b$ modulo $n$. \hyperref[chap-7]{Chapitre 7}, exercices 3, 7, 10.\par\vspace{1.7mm}
\noindent\textbf{A07} · Équations diophantiennes et contraintes sur les solutions. \hyperref[chap-7]{Chapitre 7}, exercices 4, 8, 9, 11.\par\vspace{1.7mm}
\noindent\textbf{A08} · Nombres premiers, infinitude, décomposition unique et étude de primalité. \hyperref[chap-8]{Chapitre 8}, exercices 1, 2, 5, 10.\par\vspace{1.7mm}
\noindent\textbf{A09} · Petit théorème de Fermat ; limites d’un test de primalité. \hyperref[chap-8]{Chapitre 8}, exercices 3, 8, 9.\par\vspace{1.7mm}
\noindent\textbf{A10} · Crible, décomposition algorithmique et chiffrement. \hyperref[chap-8]{Chapitre 8}, exercices 4, 6, 7, 11.\par\vspace{1.7mm}
\noindent\textbf{M01} · Matrices carrées, lignes/colonnes, opérations et dimensions. \hyperref[chap-9]{Chapitre 9}, exercices 1, 2, 3, 5.\par\vspace{1.7mm}
\noindent\textbf{M02} · Inverse en ordres $2$ et $3$, résolution de systèmes. \hyperref[chap-9]{Chapitre 9}, exercices 4, 6, 7, 9, 11.\par\vspace{1.7mm}
\noindent\textbf{M03} · Transformations géométriques représentées par matrices. \hyperref[chap-9]{Chapitre 9}, exercices 8, 10, 12.\par\vspace{1.7mm}
\noindent\textbf{M04} · Puissances en ordres $2$ et $3$, récurrences et preuve de formules. \hyperref[chap-10]{Chapitre 10}, exercices 1, 2, 5, 6, 7, 9, 12.\par\vspace{1.7mm}
\noindent\textbf{M05} · Colonnes récurrentes $U_{n+1}=AU_n+C$ et problèmes d’évolution. \hyperref[chap-10]{Chapitre 10}, exercices 3, 4, 8, 10, 11.\par\vspace{1.7mm}
\noindent\textbf{G01} · Graphes, sommets, arêtes, adjacence, ordre, degrés et graphe complet. \hyperref[chap-11]{Chapitre 11}, exercices 1, 2, 3, 7.\par\vspace{1.7mm}
\noindent\textbf{G02} · Chaînes, longueurs, connexité et modèle par graphe. \hyperref[chap-11]{Chapitre 11}, exercices 4, 8, 10.\par\vspace{1.7mm}
\noindent\textbf{G03} · Adjacence, parcours de longueur fixée, comptage par puissances et preuve. \hyperref[chap-11]{Chapitre 11}, exercices 5, 6, 9, 11, 12.\par\vspace{1.7mm}
\noindent\textbf{P01} · Markov à deux/trois états : graphe pondéré, transition et distribution initiale. \hyperref[chap-12]{Chapitre 12}, exercices 1, 2, 4, 8.\par\vspace{1.7mm}
\noindent\textbf{P02} · Interprétation de $P^n$ et distribution ligne $\pi_0P^n$ ; preuves. \hyperref[chap-12]{Chapitre 12}, exercices 2, 5, 8, 11, 12.\par\vspace{1.7mm}
\noindent\textbf{P03} · Distributions invariantes à deux/trois états et étude de probabilités. \hyperref[chap-12]{Chapitre 12}, exercices 3, 6, 7, 9, 10.\par\vspace{1.7mm}
\fiche{Démontrer et approfondir}{Une progression vers des problèmes plus ouverts}
\souscours{Démonstrations mentionnées par le programme}
Conjugués de produits, inverses, puissances et binôme : chapitre 1.
Identité du module, produit et puissance : chapitre 2.
Formule d'addition : chapitre 3.
Factorisation des différences de puissances, factorisation par une racine et majoration du nombre de racines : chapitre 4.
Détermination des racines de l'unité : chapitre 5.
Écriture du PGCD comme combinaison de Bézout et Gauss : chapitre 7.
Infinitude des premiers : chapitre 8.
Comptage des parcours par les puissances d'adjacence : chapitre 11.
Probabilités de transition à plusieurs étapes et distributions lignes : chapitre 12.
\souscours{Algorithmes et outils}
Euclide et Bézout : chapitre 7, avec fonction Python.
Crible et décomposition en premiers : chapitre 8, avec fonctions Python.
Puissances, systèmes et contrôles de produits : chapitres 9 et 10.
Calcul des distributions : chapitre 12.
La fonction native \code{pow(a, n, m)} donne une exponentiation modulaire efficace ; une calculatrice ou un logiciel de calcul matriciel peut soutenir l'expérimentation, puis la preuve justifie les formules.
\souscours{Repères historiques}
Le programme propose de relier les complexes aux équations algébriques de la Renaissance, puis à leur représentation géométrique. Les chapitres 1 à 5 suivent ce passage du calcul au plan.
L'arithmétique relie division d'Euclide, identités de Bézout et applications au chiffrement.
Les graphes, matrices et chaînes de Markov éclairent l'unité entre problèmes de réseaux, calcul et probabilités. Les dossiers fournissent des contextes de recherche et de communication orale.
\souscours{Prolongements guidés}
Les carrées de complexes, suites complexes, inégalité triangulaire, interpolation, Euler, RSA, diffusion et classement par marche aléatoire sont présentés comme des problèmes guidés. Aucune théorie générale de diagonalisation, de cryptographie ou de convergence des chaînes n'est supposée.
\repere{projet-1}{Une spirale et son centre}
\repere{projet-2}{Le pentagone régulier}
\repere{projet-3}{Un message affine}
\repere{projet-4}{Pourquoi un RSA miniature fonctionne}
\repere{projet-5}{Deux calendriers}
\repere{projet-6}{Visiter toutes les liaisons}
\repere{projet-7}{Une parabole par trois mesures}
\repere{projet-8}{Diffusion d’Ehrenfest à deux particules}
\repere{projet-9}{Un modèle de classement de pages}
\fiche{Sources et démarche éditoriale}{Recherche vérifiée le 23 octobre 2026}
\begin{itemize}
\item Ministère, \href{https://www.education.gouv.fr/bo/19/Special8/MENE1921264A.htm}{arrêté du 19 juillet 2019, BO spécial n° 8 du 25 juillet 2019}, programme de mathématiques expertes.
\item Éduscol, \href{https://eduscol.education.gouv.fr/sites/default/files/document/spe264annexe1158825pdf-84165.pdf}{annexe officielle en vigueur, 11 pages}.
\item Éduscol, \href{https://eduscol.education.gouv.fr/5817/programmes-et-ressources-en-mathematiques-voie-gt}{programmes et ressources en mathématiques, voie générale et technologique}, mise à jour juillet 2026.
\item Mission Villani-Torossian, \href{https://eduscol.education.gouv.fr/sites/default/files/document/grilles-de-positionnement-des-livres-de-mathematiques-70116.pdf}{grilles de positionnement des livres de mathématiques}, mai 2019.
\end{itemize}
Les anciens cours OptiTAB sur les complexes, l'arithmétique et les matrices ont servi de repères internes ; le texte de cette édition est restructuré pour l'option et complété selon l'annexe officielle.
Les exercices, dossiers et évaluations sont des créations originales avec leurs solutions. Ils ne sont pas présentés comme des sujets officiels de baccalauréat.
Les données et réseaux sont fictifs. Les schémas sont construits à partir des coordonnées, des arêtes et des matrices du texte ; les algorithmes et les calculs clés ont fait l'objet de contrôles indépendants.
